\documentclass[a4paper, noanswers, addpoints, 12pt]{exam}
\usepackage[T1]{fontenc}
\usepackage[utf8]{inputenc}
\usepackage{lmodern}
\usepackage{tikz}
\usetikzlibrary{positioning,arrows.meta}
\usepackage[french]{babel}
\usepackage{xcolor}
\usepackage{listings}
\usepackage{graphicx}
\usepackage{array}
\usepackage{needspace}
\usepackage[most]{tcolorbox}

\definecolor{codegreen}{rgb}{0,0.6,0}
\definecolor{codegray}{rgb}{0.5,0.5,0.5}
\definecolor{codepurple}{rgb}{0.58,0,0.82}
\definecolor{backcolour}{rgb}{0.95,0.95,0.92}

\lstdefinestyle{mystyle}{
    backgroundcolor=\color{backcolour},
    commentstyle=\color{codegreen},
    keywordstyle=\color{magenta}\bfseries,
    numberstyle=\scriptsize\color{codegray},
    stringstyle=\color{codepurple},
    basicstyle=\ttfamily\small,
    breakatwhitespace=false,
    breaklines=true,
    captionpos=b,
    keepspaces=true,
    columns=fullflexible,
    numbers=left,
    numbersep=6pt,
    xleftmargin=1.6em,
    showspaces=false,
    showstringspaces=false,
    showtabs=false,
    upquote=true,
    tabsize=4,
    escapeinside={(*@}{@*)}
}
\lstset{style=mystyle, language=Python}
\lstset{extendedchars=true,literate={é}{{\'e}}1 {è}{{\`e}}1 {ê}{{\^e}}1 {à}{{\`a}}1 {ç}{{\c c}}1 {ù}{{\`u}}1 {ï}{{\"i}}1 {É}{{\'E}}1}
% Code à compléter : interligne large pour écrire à la main dans les trous.
\lstdefinestyle{trous}{style=mystyle, basicstyle=\ttfamily\small\linespread{1.45}\selectfont}
% Console : message d'erreur ou sortie, sans numéros de ligne.
\lstdefinestyle{console}{style=mystyle, language={}, numbers=none, xleftmargin=0.3em, backgroundcolor=\color{white}, frame=single, rulecolor=\color{codegray}, basicstyle=\ttfamily\footnotesize}

% Case à cocher
\newcommand{\case}{\framebox[0.95em]{\rule{0pt}{0.55em}}}
% Cadre de réponse qui commence par l'en-tête de la fonction (signature, docstring, exemple)
\newtcolorbox{cadre}{colback=white, colframe=black, boxrule=0.5pt, arc=0mm, left=1.5mm, right=1.5mm, top=0.5mm, bottom=1mm, before skip=0.3em, after skip=0.6em}
% Zones de réponse
\setlength\dottedlinefillheight{0.85cm}
\newcommand{\trou}[1]{\makebox[#1]{\dotfill}}
\newcommand{\finligne}{\leavevmode\unskip\ \dotfill\par}
\newcommand{\rep}[1]{\par\nopagebreak\fillwithdottedlines{#1}}
\newcommand{\court}{\par\nopagebreak\noindent\makebox[\linewidth]{\dotfill}\par}
% Cadre pour écrire du code (lignes pointillées)
\newcommand{\cadrecode}[1]{\par\nopagebreak\noindent\fbox{\begin{minipage}{\dimexpr\linewidth-2\fboxsep-2\fboxrule}\fillwithdottedlines{#1}\end{minipage}}\par}
% Pile vide à compléter (sommet en haut) et file vide à compléter (tête à gauche)
\newcommand{\pilevide}{\begin{tikzpicture}
  \foreach \k in {0,...,4} {\draw (0,\k*0.8) rectangle (1.5,\k*0.8+0.8);}
  \draw (0,4) -- (0,4.4); \draw (1.5,4) -- (1.5,4.4);
  \node[anchor=west, font=\small] at (1.6,3.6) {$\leftarrow$ sommet};
  \node[font=\small] at (0.75,-0.35) {pile \texttt{p}};
\end{tikzpicture}}
\newcommand{\filevide}{\begin{tikzpicture}
  \foreach \k in {0,...,5} {\draw (\k*0.95,0) rectangle (\k*0.95+0.95,0.8);}
  \node[font=\small, anchor=south west] at (0,0.85) {tête};
  \node[font=\small, anchor=south east] at (5.7,0.85) {queue};
  \node[font=\small] at (2.85,-0.35) {file \texttt{f}};
\end{tikzpicture}}
% Deux colonnes : à gauche le code, à droite les questions qui en parlent.
\newenvironment{gauche}[1][0.5]{\par\noindent\begin{minipage}[t]{#1\linewidth}\vspace{0pt}}{\end{minipage}}
\newenvironment{droite}[1][0.47]{\hfill\begin{minipage}[t]{#1\linewidth}\vspace{0pt}\raggedright}{\end{minipage}\par}



\title{CC Term NSI - POO, Pile, File et rappels de 1re}
\author{Babash - NSI Tani Malandi}
\date{Novembre 2026 (version révisée), Durée: 2h}

\usepackage[a4paper,top=2.1cm,bottom=2.5cm,left=2cm,right=2cm,headsep=0.6cm,footskip=1.2cm]{geometry}
\pointpoints{point}{points}
\begin{document}
\runningfooter{}{\footnotesize Babash - NSI Tani Malandi \textperiodcentered{} CC BY-SA 4.0 \textperiodcentered{} page \thepage\ / \numpages}{}
\firstpagefooter{}{\footnotesize Babash - NSI Tani Malandi \textperiodcentered{} CC BY-SA 4.0 \textperiodcentered{} page \thepage\ / \numpages}{}
\pagestyle{headandfoot}
\firstpageheader{\makebox[\textwidth]{Nom : \dotfill\ Pr\'enom : \dotfill\ Classe : \makebox[3cm]{\dotfill}}}{}{}

% Titre (même rendu que \maketitle, plus serré)
\begin{center}
{\LARGE CC Term NSI - POO, Pile, File et rappels de 1re\par}\vspace{0.6em}
{\large Babash - NSI Tani Malandi\par}\vspace{0.4em}
{\large Novembre 2026 (version révisée), Durée: 2h\par}
\end{center}
\thispagestyle{headandfoot}
\begin{center}\footnotesize\itshape Exercices 1 à 5 écrits pour ce contrôle. Exercice 6 : exercice 2 du sujet 24-NSIJ2JA1 (baccalauréat 2024, Asie, jour 2), dont on garde les questions 1, 2, 4, 5, 6 et 8, reprises telles quelles ; la question 3 est remplacée par une phrase de l'énoncé, la question 7 est retirée ; les cadres de réponse donnent l'en-tête des fonctions à écrire.\end{center}

\noindent\begin{minipage}[t]{0.55\linewidth}\vspace{0pt}
\textbf{\large Consignes :}\par\smallskip
\fbox{\parbox{\dimexpr\linewidth-2\fboxsep-2\fboxrule}{\small
On répond directement sur le sujet, dans la place prévue.\\[2pt]
Chaque question demande une seule chose : une réponse courte suffit, sauf si une phrase est demandée.\\[2pt]
Les exercices sont indépendants. Les durées indiquées sont des conseils.}}
\end{minipage}\hfill
\begin{minipage}[t]{0.42\linewidth}\vspace{0pt}
\textbf{\large Barème :}\par\smallskip
\small
\begin{tabular}{|l|c|c|}\hline
\textbf{Exercice} & \textbf{Points} & \textbf{Durée} \\\hline
1. Cours & 4 & 12 min \\\hline
2. POO : lire & 2,5 & 15 min \\\hline
3. POO : compléter & 1,5 & 8 min \\\hline
4. Pile et file & 3 & 20 min \\\hline
5. Dictionnaire & 1,5 & 8 min \\\hline
6. Type bac & 7,5 & 45 min \\\hline
Relecture &  & 12 min \\\hline
\textbf{Total} & \textbf{20} & \textbf{2 h} \\\hline
\end{tabular}
\end{minipage}

% ---------------------------------------------------------------- 1
\vspace{-0.3cm}
\section{Questions de cours (4 points)}
\begin{center}\textit{Thèmes abordés : programmation orientée objet, pile et file. Durée conseillée : 12 min.}\end{center}

\begin{questions}
\question[\half] Que signifie le sigle LIFO (en anglais ou en français) ?
\court

\question[\half] Dans une méthode, que désigne le paramètre \textit{self} ? Cocher la bonne réponse.
\par\smallskip
\begin{tabular}{@{}p{0.47\linewidth}p{0.47\linewidth}@{}}
\case\ la classe où la méthode est écrite & \case\ l'objet sur lequel la méthode est appelée \\[3pt]
\case\ le premier attribut de la classe & \case\ la valeur que la méthode renvoie \\
\end{tabular}

\question Pour chaque affirmation, cocher V (vraie) ou F (fausse).
\par\smallskip
\begin{parts}
\part[\half] \parbox[t]{0.72\linewidth}{La méthode \texttt{\_\_init\_\_} est appelée automatiquement quand on crée un objet.}\hfill\case\,V\quad\case\,F
\part[\half] \parbox[t]{0.72\linewidth}{Un attribut est une fonction définie à l'intérieur d'une classe.}\hfill\case\,V\quad\case\,F
\part[\half] \parbox[t]{0.72\linewidth}{Dans une pile, l'élément retiré en premier est le dernier ajouté.}\hfill\case\,V\quad\case\,F
\part[\half] \parbox[t]{0.72\linewidth}{Dans une file, l'élément retiré en premier est le dernier ajouté.}\hfill\case\,V\quad\case\,F
\end{parts}

\question[1] Relier chaque situation à la structure qui la modélise.
\par\smallskip
\begin{center}
\begin{tikzpicture}[every node/.style={font=\normalsize}]
\node[anchor=east] (s1) at (0,1.5) {les documents envoyés à l'imprimante de la salle};
\node[anchor=east] (s2) at (0,0.75) {le bouton « Retour » d'un navigateur web};
\node[anchor=east] (s3) at (0,0) {l'annulation de la dernière action (Ctrl + Z)};
\node[anchor=east] (s4) at (0,-0.75) {les clients qui attendent à la caisse d'un magasin};
\foreach \k/\y in {1/1.5,2/0.75,3/0,4/-0.75} {\fill (0.25,\y) circle (2pt);}
\fill (2.5,1.1) circle (2pt); \node[anchor=west] at (2.65,1.1) {pile};
\fill (2.5,-0.35) circle (2pt); \node[anchor=west] at (2.65,-0.35) {file};
\end{tikzpicture}
\end{center}

% ---------------------------------------------------------------- 2
\needspace{24\baselineskip}
\section{POO : lire une classe et créer des objets (2,5 points)}
\begin{center}\textit{Thème abordé : programmation orientée objet. Durée conseillée : 15 min.}\end{center}

La barge relie Petite-Terre et Grande-Terre. On modélise une barge par la classe suivante.
\par\noindent\begin{minipage}{\linewidth}
\begin{lstlisting}[basicstyle=\ttfamily\footnotesize]
class Barge:
    def __init__(self, nom, capacite):
        self.nom = nom
        self.capacite = capacite
        self.passagers = 0

    def places_libres(self):
        return self.capacite - self.passagers

    def embarquer(self, n):
        if n <= self.places_libres():
            self.passagers = self.passagers + n
            return True
        return False
\end{lstlisting}
\end{minipage}\par

\question[\half] Donner les noms des attributs de la classe \textit{Barge}. \finligne

\question Écrire l'instruction demandée (une ligne chacune).
\begin{parts}
\part[\half] Créer une barge nommée \textit{"Corail"}, qui peut transporter 400 personnes, et la stocker dans la variable \textit{b1}.
\court
\part[\half] Faire embarquer 120 personnes dans la barge \textit{b1}.
\court
\end{parts}

\needspace{9\baselineskip}
\question On exécute le programme suivant.
\par\smallskip
\begin{gauche}[0.58]
\begin{lstlisting}[basicstyle=\ttfamily\footnotesize]
b = Barge("Lagon", 300)
print(b.embarquer(250))
print(b.embarquer(80))
print(b.places_libres())
b.passagers = b.passagers - 100
print(b.embarquer(80), b.passagers)
\end{lstlisting}
\end{gauche}
\begin{droite}[0.39]
\begin{parts}
\part[\half] Qu'affiche la ligne 4 ?\\[3pt] \trou{4.5cm}
\part[\half] Qu'affiche la ligne 6 ?\\[3pt] \trou{4.5cm}
\end{parts}
\end{droite}

% ---------------------------------------------------------------- 3
\needspace{20\baselineskip}
\section{POO : compléter une classe (1,5 point)}
\begin{center}\textit{Thème abordé : programmation orientée objet, liste de tuples. Durée conseillée : 8 min.}\end{center}

Une application de musique range les morceaux d'une playlist dans une liste de tuples \textit{(titre, duree)}, où \textit{duree} est un nombre entier de secondes. La méthode \textit{duree\_totale} renvoie la durée totale de la playlist, en secondes ; une playlist vide a une durée totale de 0.
\par\noindent\begin{minipage}{\linewidth}
\begin{lstlisting}[style=trous]
class Playlist:
    def __init__(self, nom):
        self.nom = nom
        self.morceaux = []   # liste de tuples (titre, duree)

    def ajouter(self, titre, duree):
        (*@\trou{9cm}@*)

    def duree_totale(self):
        total = (*@\trou{3cm}@*)
        for titre, duree in (*@\trou{5cm}@*):
            total = (*@\trou{6cm}@*)
        return (*@\trou{3cm}@*)
\end{lstlisting}
\end{minipage}\par

\question[\half] Compléter la ligne 7 : la méthode \textit{ajouter} ajoute le tuple \textit{(titre, duree)} à la fin de la liste \textit{morceaux}.

\question[1] Compléter les lignes 10 à 13 de la méthode \textit{duree\_totale}.

% ---------------------------------------------------------------- 4
\needspace{14\baselineskip}
\section{Pile et file (3 points)}
\begin{center}\textit{Thème abordé : structures de données linéaires : pile et file. Durée conseillée : 20 min.}\end{center}

On utilise l'interface du cours :
\begin{center}\small
\begin{tabular}{|l|l|>{\raggedright\arraybackslash}p{6.4cm}|}\hline
\textbf{Pile} & \textbf{File} & \textbf{Rôle} \\\hline
\texttt{creer\_pile\_vide()} & \texttt{creer\_file\_vide()} & crée une structure vide \\\hline
\texttt{est\_vide(p)} & \texttt{est\_vide(f)} & renvoie \texttt{True} si la structure est vide \\\hline
\texttt{empiler(p, e)} & \texttt{enfiler(f, e)} & ajoute \texttt{e} au sommet (pile) ou en queue (file) \\\hline
\texttt{depiler(p)} & \texttt{defiler(f)} & retire l'élément du sommet (pile) ou de la tête (file) et le renvoie \\\hline
\texttt{sommet(p)} & \texttt{tete(f)} & renvoie l'élément du sommet ou de la tête sans le retirer \\\hline
\texttt{taille(p)} & \texttt{taille(f)} & renvoie le nombre d'éléments \\\hline
\end{tabular}
\end{center}

\needspace{12\baselineskip}
\question On exécute le programme suivant.
\par\smallskip
\begin{gauche}[0.58]
\begin{lstlisting}[basicstyle=\ttfamily\footnotesize]
p = creer_pile_vide()
f = creer_file_vide()
for x in [4, 7, 1, 9]:
    empiler(p, x)
    enfiler(f, x)
a = depiler(p)
c = defiler(f)
empiler(p, a + c)
enfiler(f, sommet(p))
print(sommet(p), tete(f), taille(p), taille(f))
\end{lstlisting}
\end{gauche}
\begin{droite}[0.39]
\begin{parts}
\part[\half] Compléter le dessin de la pile \textit{p} ci-dessous, juste après la boucle \textit{for}.
\part[\half] Compléter le dessin de la file \textit{f} ci-dessous, juste après la boucle \textit{for}.
\part[\half] Qu'affiche la ligne 10 ?\\[4pt] \trou{4.5cm}
\end{parts}
\end{droite}
\par\medskip
\begin{center}\pilevide\hspace{1.2cm}\raisebox{1.5cm}{\filevide}\end{center}

\needspace{16\baselineskip}
\question La fonction \textit{inverser} prend en paramètre une file \textit{f} et renvoie une nouvelle file qui contient les mêmes éléments dans l'ordre inverse : celui qui était à la queue est maintenant en tête.
\begin{parts}
\part[1] Compléter les quatre trous du code de la fonction \textit{inverser}.
\par\noindent\begin{minipage}{\linewidth}
\begin{lstlisting}[style=trous]
def inverser(f):
    p = creer_pile_vide()
    while not est_vide(f):
        empiler(p, (*@\trou{4.5cm}@*))
    g = creer_file_vide()
    while (*@\trou{5cm}@*):
        enfiler(g, (*@\trou{4.5cm}@*))
    return (*@\trou{3cm}@*)
\end{lstlisting}
\end{minipage}\par
\part[\half] La file \textit{f} contient, de la tête à la queue : \textit{"Anli"}, \textit{"Inès"}, \textit{"Mathis"}, \textit{"Zaïnaba"}. On exécute \textit{g = inverser(f)}. Donner le contenu de \textit{g}, de la tête à la queue.
\court
\end{parts}

% ---------------------------------------------------------------- 5
\needspace{20\baselineskip}
\section{Rappels de 1re : dictionnaire (1,5 point)}
\begin{center}\textit{Thème abordé : dictionnaire, parcours. Durée conseillée : 8 min.}\end{center}

Un commerçant enregistre les prix au kilo (en euros) dans un dictionnaire :
\par\noindent\begin{minipage}{\linewidth}
\begin{lstlisting}
prix = {"mangues": 2.5, "bananes": 2, "manioc": 1.5, "tomates": 4}
\end{lstlisting}
\end{minipage}\par

\question[\half] Qu'affiche le programme suivant ?
\par\smallskip
\begin{gauche}[0.58]
\begin{lstlisting}[basicstyle=\ttfamily\footnotesize]
for produit in prix:
    if prix[produit] >= 2.5:
        print(produit, prix[produit])
\end{lstlisting}
\end{gauche}
\begin{droite}[0.39]
\rep{1.7cm}
\end{droite}

\needspace{20\baselineskip}
\question[1] Un panier est un dictionnaire qui associe à chaque produit sa masse en kilos, par exemple \textit{\{"mangues": 2, "tomates": 0.5\}}. La fonction \textit{cout} doit renvoyer le prix total du panier : ici, \textit{7.0}. Le programme ci-dessous s'arrête avec le message d'erreur affiché sous le code. Corriger la ligne 5.
\par\noindent\begin{minipage}{\linewidth}
\begin{lstlisting}
prix = {"mangues": 2.5, "bananes": 2, "manioc": 1.5, "tomates": 4}

def cout(panier, prix):
    total = 0
    for produit in prix:
        total = total + panier[produit] * prix[produit]
    return total

print(cout({"mangues": 2, "tomates": 0.5}, prix))
\end{lstlisting}
\begin{lstlisting}[style=console]
Traceback (most recent call last):
  File "programme.py", line 9, in <module>
    print(cout({"mangues": 2, "tomates": 0.5}, prix))
          ~~~~^^^^^^^^^^^^^^^^^^^^^^^^^^^^^^^^^^^^^^
  File "programme.py", line 6, in cout
    total = total + panier[produit] * prix[produit]
                    ~~~~~~^^^^^^^^^
KeyError: 'bananes'
\end{lstlisting}
\end{minipage}\par
\medskip
\noindent Ligne 5 corrigée : \trou{11cm}

% ---------------------------------------------------------------- Bac
\newpage
\section{Exercice Type Bac : parenthèses et pile (7,5 points)}
\begin{center}
\textit{D'après le sujet 24-NSIJ2JA1, baccalauréat 2024, Asie, jour 2, exercice 2. Durée conseillée : 45 min.}
\end{center}

\textit{Cet exercice porte sur la programmation, la programmation orientée objet et les structures de données linéaires.}

\textit{Cet exercice est composé de 2 parties indépendantes.}

Dans cet exercice, on appelle parenthèses les couples de caractères \texttt{()}, \texttt{\{\}} et \texttt{[]}. Pour chaque couple de parenthèses, la première parenthèse est appelée la parenthèse ouvrante du couple et la seconde est appelée la parenthèse fermante du couple.

On dit qu'une expression (chaine de caractères) est bien parenthésée si
\begin{itemize}
\item chaque parenthèse ouvrante correspond à une parenthèse fermante de même type ;
\item les expressions comprises entre parenthèses sont des expressions bien parenthésées.
\end{itemize}

Par exemple, l'expression \lstinline|'tab[2*(i + 4)] - tab[3]'| est bien parenthésée.

En revanche, l'expression \lstinline|'tab[2*(i + 4] - tab[3)'| n'est pas bien parenthésée, car la première parenthèse fermante \texttt{]} devrait correspondre à la dernière parenthèse ouvrante \texttt{(}.

\question[1] Déterminer si l'expression \lstinline|'[2*(i+1)-3) for i in range(3, 10)]'| est bien parenthésée. Justifier votre réponse.
%! bac
\rep{2.5cm}

\subsection*{Partie A}
On peut observer que, si les parenthèses vont par couple, une expression bien parenthésée contient autant de parenthèses ouvrantes que de parenthèses fermantes.

On se propose d'écrire une fonction qui vérifie si une chaine de caractères est bien parenthésée.

\question[1] Écrire une fonction \textit{compte\_ouvrante} qui prend en paramètre une chaine de caractères \textit{txt} et qui renvoie le nombre de parenthèses ouvrantes qu'il contient.
%! bac
\begin{cadre}
\begin{lstlisting}[numbers=none, xleftmargin=0pt, backgroundcolor=\color{white}]
def compte_ouvrante(txt):
    """Renvoie le nombre de parenthèses ouvrantes ( [ { de txt.
    Exemple : compte_ouvrante('tab[2*(i + 4)] - tab[3]') renvoie 3."""
\end{lstlisting}
\fillwithdottedlines{4.3cm}
\end{cadre}

\needspace{12\baselineskip}
\noindent\textit{On dispose aussi d'une fonction} \textit{compte\_fermante}\textit{, écrite sur le même modèle, qui prend en paramètre une chaine de caractères} \textit{txt} \textit{et qui renvoie le nombre de parenthèses fermantes qu'il contient. (Phrase ajoutée au sujet d'origine, à la place de sa question 3.)}

\question[\half] En utilisant les deux fonctions précédentes, écrire une fonction \textit{bon\_compte} qui prend en paramètre une chaine de caractères \textit{txt} et qui renvoie \textit{True} si \textit{txt} a autant de parenthèses ouvrantes que parenthèses fermantes et \textit{False} sinon.
%! bac
%! pas-fonction (une ligne à écrire : return compte_ouvrante(txt) == compte_fermante(txt))
\begin{cadre}
\begin{lstlisting}[numbers=none, xleftmargin=0pt, backgroundcolor=\color{white}]
def bon_compte(txt):
\end{lstlisting}
\fillwithdottedlines{1.7cm}
\end{cadre}

\question[\half] Donner un exemple de chaine de caractères pour laquelle \textit{bon\_compte} renvoie \textit{True} alors qu'elle n'est pas bien parenthésée.
%! bac
\court

\needspace{20\baselineskip}
\subsection*{Partie B}
Comme l'algorithme précédent n'est pas suffisant, on se propose d'implémenter un algorithme utilisant une structure linéaire de pile.

On se propose d'écrire une classe \textit{Pile} qui implémente la structure de pile.
\par\noindent\begin{minipage}{\linewidth}
\begin{lstlisting}[style=trous]
class Pile:
    def __init__(self):
        self.contenu = []

    def est_vide(self):
        return len(self.contenu) == 0

    def empiler(self, elt):
        (*@\trou{8cm}@*)

    def depiler(self):
        if self.est_vide():
            return "La pile est vide."
        return (*@\trou{7cm}@*)
\end{lstlisting}
\end{minipage}\par

\question[1\half] Compléter les lignes 9 et 14 du code précédent pour que la méthode \textit{empiler} permette d'empiler un élément \textit{elt} dans une pile et que la méthode \textit{depiler} permette de dépiler une pile en renvoyant l'élément dépilé. Vous n'écrirez que le code des deux méthodes.
%! bac
%! une-demande (deux lignes d'une même classe, complétées sur le sujet)

\needspace{32\baselineskip}
Un algorithme permettant de vérifier si une expression est bien parenthésée consiste à
\begin{itemize}
\item créer une pile vide \textit{p} ;
\item parcourir l'expression en testant chaque caractère :
  \begin{itemize}
  \item si c'est une parenthèse ouvrante, on l'empile dans \textit{p} ;
  \item si c'est une parenthèse fermante, on dépile \textit{p} ;
    \begin{itemize}
    \item si les deux caractères correspondent à un couple de parenthèses, on continue le parcours,
    \item sinon l'expression n'est pas bien parenthésée ;
    \end{itemize}
  \item sinon on continue le parcours.
  \end{itemize}
\end{itemize}
Si l'expression a été entièrement parcourue, on teste la pile ; l'expression est bien parenthésée si la pile est vide.

\needspace{14\baselineskip}
\question[3] Écrire une fonction \textit{est\_bien\_parenthesee} qui prend en paramètre une chaine de caractères et qui implémente l'algorithme précédent.
%! bac
\begin{cadre}
\begin{lstlisting}[numbers=none, xleftmargin=0pt, backgroundcolor=\color{white}]
def est_bien_parenthesee(txt):
    """Renvoie True si txt est bien parenthésée, False sinon.
    Exemples : est_bien_parenthesee('tab[2*(i + 4)] - tab[3]') renvoie True ;
               est_bien_parenthesee('(]') renvoie False."""
\end{lstlisting}
\fillwithdottedlines{7.2cm}
\end{cadre}

\end{questions}

\end{document}
